%% sec7_2.tex — Section 7.2: Annealing Strategies for Sampling
%% Chapter 7 of "A First Course in Monte Carlo Methods"
%% Sanz-Alonso & Al-Ghattas, University of Chicago

\section{Annealing Strategies for Sampling}
\label{sec:7.2}

Sampling multi-modal target distributions that have regions of high probability separated by
regions of low probability can be challenging. Annealing strategies seek to improve sampling by
introducing auxiliary \textit{tempered} distributions that bridge regions of low probability.

\subsection{Simulated Tempering}
\label{sec:7.2.1}

Let $\beta_1 = 1 > \beta_2 > \cdots > \beta_K$ be $K$ \textit{decreasing} inverse temperatures
and let $f_{\beta_k}(x) \propto f(x)^{\beta_k}$. Here $f = f_{\beta_1}$ is the target
distribution of interest and, for $1 \leq k \leq K$, the density $f_{\beta_k}$ is a
\textit{flattened} version of $f$, which is easier to sample. Simulated tempering considers an
augmented target distribution
\[
  f^{\mathrm{ST}}(x, k) \propto c_k f_{\beta_k}(x)
\]
over the variable of interest $x$ and the discrete auxiliary variable $k$ which indexes the
inverse temperatures. Samples from the distribution of interest $f$ can then be obtained by
keeping the $x$-variable of the samples $(X^{(n)}, k^{(n)})$ with $k^{(n)} = 1$.

A simple way to design a Markov chain Monte Carlo algorithm for the extended distribution
$f^{\mathrm{ST}}$ is by alternating between updates of the $x$ and $k$ components. Moves for
the latter can be proposed using the simple kernel
\begin{align*}
  r(k, k+1) = r(k, k-1) &= \tfrac{1}{2}, \quad k \in \{2, \ldots, K-1\},\\
  r(K, K-1) = r(1, 2)   &= 1.
\end{align*}
The simulated tempering algorithm then proceeds as follows:

\begin{algorithm}[H]
\caption{Simulated Tempering}
\label{alg:7.2}
\begin{algorithmic}[1]
\Require Inverse temperatures $\beta_1 = 1 > \cdots > \beta_K$, proposal Markov kernel
  $q(x, z)$, proposal Markov kernel $r(k, \ell)$, initialization $(X^{(0)}, q^{(0)})$,
  sample size $N$.
\For{$n = 0, 1, \ldots, N-1$}
  \State Draw $\ell^* \sim r(k^{(n)}, \cdot)$.
  \State Update
    \[
      k^{(n+1)} =
      \begin{cases}
        \ell^*   & \text{with probability } a_{k,\ell} := \min\!\left\{1,\,
                   \dfrac{f^{\mathrm{ST}}(X^{(n)}, \ell^*)}{f^{\mathrm{ST}}(X^{(n)}, k^{(n)})}
                   \cdot \dfrac{r(k^{(n)}, \ell^*)}{r(\ell^*, k^{(n)})}\right\}, \\[6pt]
        k^{(n)} & \text{with probability } 1 - a_{k,\ell}.
      \end{cases}
    \]
  \State Draw $Z^* \sim q(X^{(n)}, \cdot)$.
  \State Update
    \[
      X^{(n+1)} =
      \begin{cases}
        Z^*      & \text{with probability } a_{X,Z} := \min\!\left\{1,\,
                   \dfrac{f^{\mathrm{ST}}(Z^*, k^{(n+1)})}{f^{\mathrm{ST}}(X^{(n)}, k^{(n+1)})}
                   \cdot \dfrac{q(Z^*, X^{(n)})}{q(X^{(n)}, Z^*)}\right\}, \\[6pt]
        X^{(n)} & \text{with probability } 1 - a_{X,Z}.
      \end{cases}
    \]
\EndFor
\Ensure Sample $\{(X^{(n)}, k^{(n)})\}_{n=1}^N$.
\end{algorithmic}
\end{algorithm}

The idea is that for low inverse temperature the chain will easily move across the state space,
and by the time the augmented chain returns to the temperature of interest $\beta_1 = 1$ the
$x$ variable will have moved to a different region of the state space, accelerating the mixing.

\begin{center}
\fbox{\parbox{0.92\textwidth}{%
\textbf{Pros and Cons}: Simulated tempering faces two implementation issues. First, how should
we choose the constants $c_k$? It is generally recommended that the temperatures are chosen so
that the extended chain spends roughly the same amount of time at each temperature. If all the
$f_{\beta_k}$ can be normalized, then the constants $c_k$ should then be chosen equal. However,
in practice these normalizing constants are typically unknown and need to be estimated using an
alternative adaptive algorithm (e.g.\ Wang--Landau). The second implementation issue is how to
choose the inverse temperatures. There is a trade-off between having sufficient acceptance
probability for moving from one temperature to another, and not having too many temperature
levels. There is a vast literature on this issue, notably in physics.
}}
\end{center}

\subsection{Parallel Tempering}
\label{sec:7.2.2}

Parallel tempering also uses a decreasing sequence of inverse temperatures $\beta_k$. However,
it runs $K$ chains in parallel such that the $k$-th chain has stationary distribution
$f_{\beta_k}$. The chains are allowed to communicate by a swapping mechanism, so that chains
with larger inverse temperatures benefit from the enhanced mixing of the chains with lower
inverse temperature. Precisely, after updating each of the $K$ chains, we pick a pair of
chains and propose to swap their states. This proposed swapping is followed by an accept/reject
step to retain the correct invariant distribution. As in simulated tempering, we set
$\beta_1 = 1$ so that $f_{\beta_1} = f$. We denote by $X_k^{(n)}$ the $n$-th sample from
the $k$-th chain.

\begin{algorithm}[H]
\caption{Parallel Tempering}
\label{alg:7.3}
\begin{algorithmic}[1]
\Require Inverse temperatures $\beta_1 = 1 > \cdots > \beta_K$, proposal Markov kernels
  $\{q_k(x, z)\}_{k=1}^K$, initializations $\{X_k^{(0)}\}_{k=1}^K$, sample size $N$.
\For{$n = 0, 1, \ldots, N-1$}
  \State For $k = 1, \ldots, K$: generate $\tilde{X}_k^{(n+1)}$ by doing a Metropolis
    Hastings step (including accept/reject) with current state $X_k^{(n)}$, proposal kernel
    $q_k$, and target $f_{\beta_k}$.
  \State Choose $\ell, m \in \{1, \ldots, K\}$ with $\ell \neq m$ uniformly at random.
  \State For $k \notin \{\ell, m\}$ set $X_k^{(n+1)} = \tilde{X}_k^{(n+1)}$.
  \State Attempt a swap of states between the $\ell$-th and the $m$-th chains:
    \[
      \bigl(X_\ell^{(n+1)}, X_m^{(n+1)}\bigr) =
      \begin{cases}
        \bigl(\tilde{X}_m^{(n+1)},\, \tilde{X}_\ell^{(n+1)}\bigr)
          & \text{with probability } a_{\ell,m},\\[4pt]
        \bigl(\tilde{X}_\ell^{(n+1)},\, \tilde{X}_m^{(n+1)}\bigr)
          & \text{with probability } 1 - a_{\ell,m},
      \end{cases}
    \]
    where
    \[
      a_{\ell,m} := \min\!\left\{1,\;
        \frac{f_{\beta_\ell}\!\bigl(\tilde{X}_m^{(n+1)}\bigr)\,
              f_{\beta_m}\!\bigl(\tilde{X}_\ell^{(n+1)}\bigr)}
             {f_{\beta_m}\!\bigl(\tilde{X}_m^{(n+1)}\bigr)\,
              f_{\beta_\ell}\!\bigl(\tilde{X}_\ell^{(n+1)}\bigr)}\right\}.
    \]
\EndFor
\Ensure Sample $= \{X_1^{(n)}\}_{n=1}^N$.
\end{algorithmic}
\end{algorithm}

Parallel tempering requires proposal Markov kernels $\{q_k\}_{k=1}^K$ to update each chain.
A popular choice is to use the Langevin proposals described in Chapter~6. Precisely, for the
$k$-th chain one can define a proposal Markov kernel $q_k$ by proposing moves
$X_k^{(n)} \mapsto Z_k^*$ according to
\[
  Z_k^* = X_k^{(n)} + \epsilon_k \nabla\log f\!\bigl(X_k^{(n)}\bigr)
          + \sqrt{2\epsilon_k \beta_k^{-1}}\, \xi_k^*,
\]
where $\xi_k^* \sim \Normal(0, I)$ and $\epsilon_k > 0$ is a given step size. Notice that large
$\beta_k^{-1}$ results in large random moves that can help explore the state space. Such moves
may still be accepted with high probability, since they are targeting the flattened distribution
$f_{\beta_k}$.

\begin{center}
\fbox{\parbox{0.92\textwidth}{%
\textbf{Pros and Cons}: In contrast to simulated tempering, parallel tempering does not require
estimating the normalizing constants $c_k$. While parallel tempering involves running $K$
chains in parallel, doing so is not a significant obstacle with modern parallel computing. The
question of how to choose the number $K$ of chains and the temperatures $\beta_k$ has been
widely studied, and there is some evidence that geometric scaling of the temperatures may be
optimal under mild assumptions.
}}
\end{center}
